\section{Position du problème}

\subsection{D'un courtier en énergie à un organisme de formation}

Audit Énergie s'est d'abord construite autour du courtage en énergie. Son métier consiste
à accompagner des clients professionnels dans la compréhension, la comparaison et la
renégociation de leurs contrats, sur un marché devenu progressivement plus complexe. Au
fil des années, cette activité a forgé une expertise concrète des métiers de la vente à
distance : la prospection, la relation client par téléphone, l'analyse d'un besoin
commercial, l'argumentation et le suivi d'un dossier dans la durée.

Ce positionnement initial éclaire le sujet du mémoire, mais il n'en constitue pas le cœur.
Le projet étudié s'inscrit dans une évolution plus récente : la transformation d'Audit
Énergie en organisme de formation. Cette évolution s'est faite en deux temps. D'abord en
interne, pour former ses propres téléprospecteurs aux fondamentaux de la relation client
et aux spécificités du secteur. Puis, depuis environ un an, de manière beaucoup plus
structurée, avec le développement d'une offre de formation externe. L'entreprise prépare
désormais des apprenants à des titres professionnels reconnus par l'État, notamment le
titre de Conseiller Relation Client à Distance ou celui d'Employé Commercial.

Ce changement d'échelle modifie en profondeur les contraintes de l'entreprise. Former
quelques collaborateurs en interne et faire fonctionner un organisme de formation externe
ne relèvent pas de la même organisation. Dans le premier cas, la transmission peut
s'appuyer sur l'expérience des salariés, des échanges directs et des supports peu
formalisés. Dans le second, il faut produire des parcours reproductibles, préparer des
contenus, assurer une continuité pédagogique, suivre la présence des élèves, proposer des
exercices, répondre aux questions et garantir une qualité suffisante pour que la formation
soit réellement exploitable.

Le sujet du mémoire naît précisément de cette transition. Il ne s'agit pas seulement
d'ajouter une fonctionnalité technique à une plateforme, mais de répondre à une question
d'organisation plus large : comment une petite structure peut-elle développer une activité
de formation à distance sans disposer des moyens humains et financiers d'un organisme de
formation déjà établi ?

\subsection{La formation à distance : une opportunité freinée par son coût}

La formation à distance représente, sur le papier, une opportunité idéale pour une petite
entreprise. Elle permet de toucher des apprenants sans salle physique, de mutualiser des
contenus, de planifier des sessions souples et de réutiliser des supports d'une promotion à
l'autre. Elle semble donc parfaitement compatible avec les contraintes d'une structure
légère.

Cette impression est pourtant incomplète. Une formation à distance n'est rentable que si
ses coûts de production et d'animation restent maîtrisés. Or elle ne se résume pas à
déposer des documents en ligne : il faut concevoir un déroulé pédagogique, produire les
cours, les rendre accessibles, accompagner les apprenants, suivre leur participation et
vérifier leur compréhension. Ces tâches sont répétitives, mais indispensables.

Le modèle le plus naturel aurait été de faire appel à des formateurs freelance pour
assurer les cours en ligne. Ce modèle a des avantages évidents : un formateur humain
adapte son discours, répond aux questions, reformule une notion difficile, donne des
exemples et entretient une dynamique de groupe. Surtout, il incarne une responsabilité
pédagogique immédiate : si un élève ne comprend pas, il intervient sur-le-champ.

Mais ce modèle se heurte à un obstacle économique fort. Pour une petite structure, le coût
récurrent d'un formateur absorbe une part importante de la valeur produite par la
formation, et cette part croît avec le nombre d'heures dispensées. Dans un marché où les
marges doivent financer l'organisation, le suivi administratif, les outils et le
développement commercial, dépendre fortement d'un intervenant externe rend le modèle
difficile à industrialiser. Le problème initial peut donc se formuler simplement :
l'entreprise voulait développer une activité de formation externe, mais le modèle
classique du formateur en direct n'était compatible ni avec ses objectifs de coût, ni avec
son besoin de passage à l'échelle.

\subsection{Première réponse : séparer la production du cours de sa diffusion}

La première réponse a consisté à ne plus raisonner en cours en direct, mais en cours
préparé à l'avance. L'idée était de construire une plateforme capable d'héberger des
contenus audio et de les diffuser automatiquement le jour prévu de la formation. Plutôt
que de rémunérer un formateur pour animer plusieurs heures de cours, le travail était
scindé en deux étapes : la préparation du contenu, puis sa lecture sous forme d'audio.

Concrètement, les cours étaient rédigés pendant la semaine, éventuellement avec l'aide
d'outils d'intelligence artificielle, puis enregistrés et déposés dans la plateforme. Le jour
de la session, la diffusion se déclenchait selon un calendrier horodaté. Pour l'élève, le
cours se présentait comme une séance organisée dans le temps ; pour l'entreprise, la
production devenait plus facilement planifiable qu'un cours en direct. En répartissant
l'effort entre une personne qui écrit et une personne qui lit, le coût total devenait
nettement inférieur à celui d'un formateur freelance.

Cette organisation a validé plusieurs hypothèses importantes : les apprenants pouvaient
suivre une formation à partir de contenus audio diffusés à distance, l'entreprise pouvait
internaliser une partie de la production pédagogique, et la plateforme devenait l'outil
central d'orchestration de la formation, et non un simple espace de dépôt.

Elle ne résolvait toutefois qu'une partie du problème. La dépendance au formateur en
direct diminuait, mais la dépendance au travail humain demeurait. Les cours devaient
toujours être rédigés, relus, découpés, enregistrés, déposés et vérifiés. À chaque nouvelle
formation, voire à chaque mise à jour, une chaîne de production manuelle devait être
relancée. Le gain existait, mais il restait limité par le caractère artisanal du processus :
l'entreprise avait déplacé le coût humain plus qu'elle ne l'avait supprimé.

\subsection{De l'IA d'assistance à l'ambition d'une plateforme autonome}

L'arrivée des outils d'intelligence artificielle générative a déplacé la manière même de
poser le problème. Dans un premier temps, l'IA pouvait servir d'aide à la rédaction : un
collaborateur produisait une première version d'un cours, reformulait un passage,
structurait une séquence. L'humain gardait le contrôle et utilisait l'IA pour accélérer
certaines tâches.

Mais le contexte technologique a vite ouvert une hypothèse plus ambitieuse. Si un modèle
de langage peut rédiger un contenu, si un outil de synthèse vocale peut le lire, si un
système peut générer des supports visuels et si un assistant peut répondre aux questions
des élèves, alors il devient envisageable de concevoir une plateforme bien plus autonome.
Dans cette perspective, l'IA cesse d'être une aide ponctuelle pour devenir une composante
centrale de la chaîne de production et d'accompagnement.

Cet objectif doit être formulé sans détour : il s'agissait bien de remplacer une part
importante des interventions humaines initialement nécessaires à la formation. Dans une
petite entreprise, ce choix n'est pas seulement technologique, il est économique et
stratégique. L'enjeu est de rendre possible une activité qui serait trop coûteuse si elle
reposait entièrement sur des ressources humaines mobilisées en continu.

Remplacer une fonction humaine ne revient cependant pas à remplacer une personne par
un appel d'API. Un formateur ne se contente pas de lire un texte : il structure un cours,
choisit des exemples, respecte un programme, s'adapte au niveau des élèves, répond aux
questions, évalue la compréhension et maintient une cohérence entre les supports.
Automatiser cette fonction suppose donc d'automatiser une chaîne complète de tâches, et
non un geste isolé. C'est cette prise de conscience qui donne au sujet sa véritable
dimension d'ingénierie. La question n'est plus « quel outil d'IA utiliser ? », mais «
comment organiser une chaîne logicielle capable de produire, contrôler, diffuser et
exploiter des contenus pédagogiques générés par IA dans un contexte réel de formation ? »

\subsection{Les essais de synthèse vocale et leurs limites}

La première brique d'automatisation envisagée après la rédaction assistée a été la synthèse
vocale. L'idée était directe : si le cours est écrit à l'avance, une voix artificielle peut le
lire. Les outils récents produisent des voix bien plus naturelles que les anciens systèmes
de lecture automatique, ce qui rend la solution séduisante pour une plateforme audio.

Les premiers essais ont montré que le problème était plus subtil que prévu. Produire
plusieurs heures d'audio n'a rien d'une courte démonstration vocale. Les textes devaient
être découpés en portions compatibles avec l'outil, envoyés successivement, puis les
fragments audio obtenus devaient être réassemblés pour reconstituer un cours continu. Ce
procédé cumulait trois difficultés. Le temps de manipulation, d'abord : même générée par
IA, la voix exigeait un découpage et un assemblage largement manuels. Le coût, ensuite :
les meilleurs outils deviennent onéreux dès qu'on produit des journées complètes de
formation. La qualité pédagogique, enfin : une voix naturelle ne suffit pas si le texte lu
n'est pas pensé pour l'oral, s'il est répétitif ou s'il ne respecte pas le rythme d'un cours.

Ces essais ont clarifié un point essentiel : la synthèse vocale n'automatise que la lecture,
pas la conception du cours. Or un audio n'est exploitable que si le texte d'origine est clair,
progressif, calibré en durée, fidèle au programme et rédigé dans un style oral. La synthèse
vocale a donc déplacé le centre du problème vers l'amont. La vraie question n'était plus de
savoir comment générer une voix, mais comment générer un script pédagogique
suffisamment fiable pour être confié à une voix automatique et diffusé à de vrais élèves.

\subsection{Le cours préenregistré ne remplace pas le formateur}

Même avec un texte et une voix satisfaisants, le cours audio préenregistré bute sur une
limite structurelle : il ne peut pas interagir. Dans un cours en direct, l'élève pose une
question, demande une reformulation, signale une incompréhension, obtient un exemple
supplémentaire. Cette interaction fait partie intégrante de la valeur d'un formateur.

Dans un modèle purement audio, elle disparaît. Les élèves écoutent, mais ne disposent
d'aucun interlocuteur pour répondre à leurs questions. La limite est d'autant plus sensible
dans une formation professionnalisante, où les profils sont variés et les besoins de
clarification fréquents : un passage limpide pour l'un peut rester obscur pour l'autre.

Cette difficulté révèle une seconde dimension du problème. Il ne suffit pas d'automatiser
la production des cours ; il faut aussi automatiser une partie de l'accompagnement. Sans
cela, la plateforme reste un simple diffuseur d'audios, capable de réduire des coûts mais
incapable de remplacer les fonctions d'échange et de clarification d'un formateur. C'est
précisément le rôle d'un assistant conversationnel fondé sur la récupération d'information
(RAG) : non pas discuter librement comme un chatbot généraliste, mais répondre aux
questions des élèves à partir du contenu réel de la formation. La distinction est
fondamentale. Un assistant pédagogique ne doit pas produire une réponse simplement
plausible ; il doit s'appuyer sur les cours et les supports disponibles. Il doit être
contextualisé.

\subsection{Une classe de problèmes au-delà d'Audit Énergie}

Le cas d'Audit Énergie est singulier par son histoire et ses contraintes, mais le problème
qu'il pose dépasse largement cette entreprise. De nombreuses petites structures connaissent
aujourd'hui la même tension : elles veulent proposer des services plus riches, plus
personnalisés et plus accessibles, sans disposer des moyens humains pour les produire
manuellement. L'intelligence artificielle générative apparaît alors comme une issue : elle
sait rédiger, reformuler, résumer, produire de l'audio, générer des supports ou répondre à
des questions.

Mais son intégration dans un contexte professionnel ne peut reposer sur des usages
ponctuels. Un collaborateur qui demande à un modèle de rédiger un cours obtient parfois
un résultat utile, sans qu'il soit pour autant traçable, reproductible, contrôlé ni
directement exploitable en production. La classe de problèmes étudiée peut donc se
formuler ainsi : comment transformer des capacités d'IA générative, le plus souvent
utilisées de façon ponctuelle, en une chaîne logicielle fiable, auditable et intégrée à un
processus métier réel ?

Dans le cas de la formation, cette question prend une forme particulièrement exigeante. Il
ne s'agit pas seulement de produire du texte, mais un objet pédagogique complet : un cours
qui respecte un programme, s'adapte au niveau des apprenants, reste compréhensible,
tient dans une durée, peut être transformé en audio, donne lieu à des supports visuels,
permet une évaluation et s'inscrit dans un parcours plus large. L'IA générative peut ici
donner une illusion de simplicité : un premier résultat convaincant masque les difficultés
réelles de répétabilité, de contrôle qualité, de coût et d'évaluation. Pour un prototype, il
suffit qu'un texte semble correct ; pour une plateforme utilisée avec de vrais élèves, il faut
pouvoir expliquer comment le contenu a été produit, quelles règles ont été appliquées et
comment sa qualité est mesurée.

\subsection{Les exigences qui encadrent le problème}

Trois exigences encadrent ce problème et orientent toute la suite du mémoire.

La première est la \emph{qualité pédagogique}. Dans une formation professionnelle, un
contenu n'est pas jugé sur sa fluidité ou son apparence, mais sur sa capacité à faire
progresser l'apprenant vers un objectif précis. Or les modèles de langage produisent des
textes très convaincants, qui peuvent pourtant contenir des imprécisions, des répétitions,
des exemples mal choisis ou des informations non vérifiées. Le risque est concret : un
élève peut retenir une explication inexacte ou perdre du temps sur un contenu hors
programme. Cet enjeu se double d'un problème de \emph{densité}. Lorsqu'on demande à
la chaîne de produire un long volume de cours à partir d'une source documentaire limitée,
le modèle, contraint de remplir le temps, tend à broder, à répéter et à digresser. La qualité
ne peut donc pas être traitée par une simple relecture finale : elle doit être intégrée à la
chaîne de production elle-même.

La deuxième exigence est la \emph{traçabilité}. Lorsqu'un humain produit un cours, il
peut expliquer ses choix et justifier sa progression. Lorsqu'une IA génère un contenu, cette
justification n'existe pas spontanément. Si la plateforme se contente de stocker le texte
final, il devient très difficile de savoir d'où vient une erreur : du programme de départ, du
plan généré, de la rédaction d'une section, d'une étape de réécriture ou de la calibration
audio. Sans artefacts intermédiaires, ni la correction, ni la confiance, ni l'amélioration
continue ne sont possibles. La traçabilité n'est pas un confort : c'est une condition de
l'évaluation.

La troisième exigence est d'ordre \emph{économique et organisationnel}. L'objectif n'est
pas de construire une plateforme techniquement impressionnante, mais de rendre viable un
modèle de formation compatible avec les moyens d'une petite entreprise. L'automatisation
ne supprime pas tous les coûts : les services d'IA ont un prix, les erreurs doivent être
contrôlées, les prompts maintenus, la plateforme supervisée et les contenus validés.
Évaluer le projet ne reviendra donc pas à comparer un formateur humain à un modèle
d'IA, mais à comparer deux organisations : l'une centrée sur l'intervention humaine
continue, l'autre sur une plateforme automatisée supervisée.

\subsection{Formulation de la problématique}

Le problème posé par ce mémoire se lit à deux niveaux. Au niveau d'Audit Énergie,
l'entreprise veut développer une formation externe avec des moyens humains limités : le
formateur en direct coûte trop cher, le cours préenregistré appauvrit l'interaction
pédagogique. Elle a donc besoin d'une plateforme capable d'automatiser la production, la
diffusion et une partie de l'accompagnement. À un niveau plus général, le même problème
concerne toutes les petites structures qui cherchent à industrialiser un service grâce à
l'IA : l'enjeu n'est pas tant d'utiliser ces outils que de les encadrer dans une architecture
permettant un usage fiable en production.

La problématique peut alors être formulée ainsi :

\begin{quote}
\emph{Comment concevoir une plateforme de formation autonome alimentée par
l'intelligence artificielle, capable de remplacer une part significative des interventions
humaines nécessaires à la production, à la diffusion et à l'accompagnement des cours,
tout en garantissant la qualité pédagogique, la traçabilité des contenus et la pertinence
des réponses apportées aux apprenants ?}
\end{quote}

Cette formulation met en évidence la tension centrale du projet. D'un côté, l'entreprise
cherche à automatiser fortement son modèle pour le rendre économiquement viable. De
l'autre, cette automatisation ne peut se faire au prix d'une perte de qualité, d'une absence
de contrôle ou d'une expérience dégradée pour les élèves. Le travail d'ingénierie consiste
donc à concilier autonomie et maîtrise.

\subsection{Objectifs de la solution attendue}

La problématique se décline en objectifs concrets, qui structureront l'étude de la solution :

\begin{itemize}
\item \textbf{Automatiser la génération des cours} à partir d'un cadre pédagogique donné,
en respectant une progression et un niveau attendus. Cette génération ne doit pas être une
rédaction libre, mais un processus structuré et contrôlable.
\item \textbf{Automatiser la transformation audio}, en calibrant les contenus pour la
lecture orale et en évitant tout montage manuel lourd, là où chaque fichier exigeait
auparavant une manipulation individuelle.
\item \textbf{Produire des supports complémentaires} (slides, exercices) cohérents avec le
cours, et non ajoutés après coup comme des éléments indépendants.
\item \textbf{Accompagner les apprenants} via un assistant conversationnel qui répond à
partir du contenu réel de la formation, plutôt que d'une connaissance générale non
contrôlée.
\item \textbf{Rendre le système auditable}, en associant à chaque contenu ses traces :
plan, sections, corrections, sources et indicateurs.
\item \textbf{Définir des indicateurs d'évaluation} : temps de production, nombre
d'interventions humaines, coût, respect des durées, qualité des contenus, résultats aux
exercices et pertinence des réponses de l'assistant.
\end{itemize}

Ces objectifs ne sont pas seulement des résultats à présenter en fin de mémoire : ils
orientent la conception. Pour mesurer le respect d'un plan, il faut un plan explicite ; pour
mesurer les corrections, il faut conserver les versions intermédiaires ; pour mesurer la
pertinence du RAG, il faut pouvoir identifier les sources utilisées. L'évaluation
influence donc directement l'architecture.

\subsection{Motivation, périmètre et transition}

Ce sujet présente un intérêt particulier pour un mémoire d'ingénieur, car il se situe à la
frontière d'un besoin métier concret et de problématiques techniques actuelles. Il ne s'agit
pas d'étudier l'intelligence artificielle de manière abstraite, mais de l'intégrer dans un
système réellement utilisé, avec de vrais élèves, des contraintes de coût et des exigences de
qualité. Sa valeur tient aussi à sa progression : la plateforme n'est pas née sous sa forme
actuelle, elle a évolué par étapes — diffusion d'audios, préparation manuelle, essais de
synthèse vocale, puis structuration d'une pipeline de génération et ajout d'un assistant
RAG. Chaque étape a fait surgir une nouvelle limite et obligé à reformuler le problème.
Prendre du recul sur cette trajectoire permet d'expliquer non seulement ce qui a été
développé, mais pourquoi ces choix ont été faits et quels compromis ont été acceptés.

Le périmètre du mémoire reste délimité. Il porte sur la conception d'une plateforme
autonome pour la production, la diffusion et l'accompagnement de cours à distance ; les
aspects administratifs, commerciaux ou réglementaires de l'organisme de formation ne sont
abordés que lorsqu'ils éclairent les contraintes du projet. Le travail ne prétend pas non
plus démontrer que l'IA peut remplacer toute forme d'enseignement humain : la
supervision humaine demeure nécessaire, mais elle change de rôle, davantage tournée vers
la validation, le pilotage et le traitement des cas qui dépassent le cadre prévu.

Cette position du problème fait apparaître plusieurs familles de solutions à examiner :
les plateformes de formation en ligne, qui diffusent et suivent mais ne produisent pas les
cours ; les outils de synthèse vocale, qui automatisent la lecture mais pas la conception ;
les modèles de langage, qui génèrent du texte mais doivent être encadrés ; les systèmes
RAG, qui contextualisent les réponses mais dépendent de la qualité de la recherche
documentaire. L'état de l'art, présenté dans la partie suivante, analysera ces approches
pour mieux situer, ensuite, la solution développée et en montrer l'originalité.
